\subsection{Components of Emotional Intelligence}
\label{sec:2.3.1}
Emotional intelligence spans the understanding and management of one's own emotions and the capacity to identify and respond appropriately to the emotions of others. Built into an AI system it is usually broken into four components: reading emotional cues, verbal and non-verbal, from expression, vocal modulation and demeanor; calibrating the system's own output to the climate of an interaction rather than escalating it; understanding another's state well enough to act on it and not merely to label it; and sustaining trust across repeated interactions through the ordinary social competencies of listening, clear communication and de-escalation.

What “emotion” is

That four-part inventory presupposes something the underlying science has not settled: that emotions are the kind of thing a system reads off a signal, the way a thermometer reads temperature. That premise has a serious rival.

The case for discrete, evolved emotions rests on some of the most influential work in twentieth-century psychology. Paul Ekman traveled to the Fore people of Papua New Guinea, who had almost no exposure to Western media, told them short stories, and asked them to match each to a photograph of a Western face; the Fore sorted the faces much as Westerners did, which Ekman took as evidence that facial expressions of emotion are universal and evolved rather than culturally learned \autocite{ekman1969pancultural}. Jaak Panksepp's work on subcortical stimulation in animals complemented it, locating seven primal emotional systems in structures shared across mammals and older than cortex \autocite{panksepp1998affective}. On this view fear and disgust really do feel like different things because they are different things, with different circuitry.

The rival camp went looking for that circuitry and did not find it. A large meta-analysis of brain-imaging studies, led by the psychologist Lisa Feldman Barrett, found no dedicated fear center, no dedicated disgust center — the amygdala lit up for novelty and salience of many kinds, the insula for almost every feeling there is, the same regions reappearing across different named emotions in shifting combinations \autocite{lindquist2012brain}. From this, Barrett built the theory of constructed emotion: the brain receives ambiguous interoceptive signal — a racing heart could mean fear, anger, excitement, or oncoming illness — and constructs the specific emotion as its best prediction of what that signal means, given context and a lifetime of experience with situations like it. Fear is not detected. It is a prediction the brain settles on. Later cross-cultural work reopened Ekman's own evidence: when researchers dropped the forced-choice method and let people sort freely, the Himba of Namibia did not group faces into the Western basic-emotion categories, and the Trobriand Islanders read the wide-eyed, gasping face Westerners call fear as a threat display instead \autocite{gendron2014perceptions}.

The field has not resolved this and is not close to resolving it. What is no longer seriously disputed is narrower and still matters here: there is no tidy map from named emotions to dedicated brain regions. An AI system's “emotional intelligence,” whichever of these theories turns out closer to right, is being built on an open scientific question about what it is modeling, not a settled inventory waiting to be implemented. Section~\ref{sec:9.8} follows Ekman's evidence forward into the applied form of the same dispute, where the question is no longer what emotion is but whether a face reveals it well enough to build a product on.
